\documentclass{article}
\usepackage{graphicx} % Required for inserting images
\usepackage{amsmath}
\usepackage{subcaption}

\title{coupled GL}
\author{Viktor Medvedev}
\date{March 2026}

\begin{document}

\section{Introduction}
Pulse propagation in erbium-doped fiber (EDF) can be described by the following coupled Ginzburg-Landau equations:

\begin{equation} \label{eq1}
    \begin{cases}
        \displaystyle
        \frac{\partial{u}}{\partial{z}} = \frac{i \Delta \beta}{2}u - \delta \frac{\partial{u}}{\partial{t}} - \frac{i\beta_{2}}{2} \frac{\partial^2 u}{\partial t^2} + \frac{g}{2}u + \frac{g}{2\Omega^2_{g}}\frac{\partial^2 u}{\partial t^2} + i\gamma\left(\left|u\right|^2 + \frac{2}{3} \left|v\right|^2\right)u + \frac{i\gamma}{3}v^2 u^*, \\
        \displaystyle
        \frac{\partial{v}}{\partial{z}} = \frac{i \Delta \beta}{2}v - \delta \frac{\partial{v}}{\partial{t}} - \frac{i\beta_{2}}{2} \frac{\partial^2 v}{\partial t^2} + \frac{g}{2}v + \frac{g}{2\Omega^2_{g}}\frac{\partial^2 v}{\partial t^2} + i\gamma\left(\left|v\right|^2 + \frac{2}{3} \left|u\right|^2\right)v + \frac{i\gamma}{3}u^2 v^*,
    \end{cases}
\end{equation}
where \(u\) and \(v\) are slowly varying pulse envelopes.

\section{Analytical analysis}
Let us consider first equation of Eq. \ref{eq1}, assuming \(v=0\):
\begin{equation} \label{eq2}
    \frac{\partial{u}}{\partial{z}} = \frac{i \Delta \beta}{2}u - \delta \frac{\partial{u}}{\partial{t}} - \frac{i\beta_{2}}{2} \frac{\partial^2 u}{\partial t^2} + \frac{g}{2}u + \frac{g}{2\Omega^2_{g}}\frac{\partial^2 u}{\partial t^2} + i\gamma\left|u\right|^2 u.
\end{equation}

Using traveling wave reduction in the form:
\begin{equation} \label{eq3}
    u\left(z, t\right) = y(\mathrm{w}) e^{i\left(k z-\omega t\right)}, \quad \text{where} \quad \mathrm{w} = z - c_{0} t,
\end{equation}
and evaluating corresponding derivatives:
\begin{equation} \label{eq4}
    \begin{cases}
        \displaystyle
        \frac{\partial u}{\partial z} = \left(y_{\mathrm{w}} + i k y\right) e^{i\left(k z-\omega t\right)},\\
        \displaystyle
        \frac{\partial u}{\partial t} = \left(-c_{0} y_{\mathrm{w}} - i \omega y\right) e^{i\left(k z-\omega t\right)}, \\
        \displaystyle
        \frac{\partial^2 u}{\partial t^2} = \left(c^{2}_{0}y_{\mathrm{ww}} + 2 i\omega c_{0} y_{\mathrm{w}} - \omega^{2}y\right) e^{i\left(k z-\omega t\right)},
    \end{cases}
\end{equation}
we can rewrite Eq. \ref{eq2} as:
\begin{equation} \label{eq5}
    \begin{split}
        &\left[\frac{c^{2}_{0} g}{2 \Omega^2_{g}} -i \left(\frac{c^2_{0}\beta_{2}} {2}\right)\right] \cdot y_{\mathrm{ww}} +
        \left[\delta c_{0} - 1 + \beta_{2} \omega c_{0} + i \left(\frac{c_{0} g \omega}{\Omega^2_{g}}\right)\right] \cdot y_{\mathrm{w}} +\\
        & + \left[\frac{g}{2}\left(1-\frac{\omega^{2}}{\Omega^2_{g}}\right) + i \left(\frac{\Delta \beta}{2}+\delta \omega + \frac{\beta_{2}}{2} \omega^{2} -k\right)\right]\cdot y +
        i \gamma y^{3} = 0.
    \end{split}
\end{equation}

Extracting real and imaginary parts of  Eq. \ref{eq5}, we obtain over-defined system of equations in the form:
\begin{equation} \label{eq6}
    \begin{cases}
        \displaystyle \left(\frac{c^{2}_{0} g}{2 \Omega^2_{g}}\right) \cdot y_{\mathrm{ww}} + \left(\delta c_{0} - 1 + \beta_{2} \omega c_{0}\right) \cdot y_{\mathrm{w}} + \left[\frac{g}{2}\left(1-\frac{\omega^{2}}{\Omega^2_{g}}\right)\right] \cdot y = 0,\\
        \displaystyle \left(-\frac{\beta_{2} c^{2}_{0}}{2}\right) \cdot y_{\mathrm{ww}} + \left(\frac{c_{0} g \omega}{\Omega^2_{g}}\right) \cdot y_{\mathrm{w}} + \left(\frac{\Delta \beta}{2} + \delta \omega + \frac{\beta_{2} \omega^{2}}{2} - k\right) \cdot y + \gamma \cdot y^3 = 0.
    \end{cases}
\end{equation}

At parameters:
\begin{equation} \label{eq7}
    \begin{cases}
    \displaystyle \delta = \frac{1-\beta_{2} c_{0} \omega}{c_{0}},\\
    \displaystyle g = 0,
    \end{cases}
\end{equation}
the equation for real part (first equation from system \ref{eq6}) can be satisfied, while imaginary part equation can be written as:
\begin{equation} \label{eq8}
    -\left(\frac{\beta_{2} c_{0}^{2}}{2} \right)y_{\mathrm{ww}} + \left(-\frac{\beta_{2} \omega^{2}}{2}+\frac{\omega}{c_{0}}-k +\frac{\Delta \beta}{2}\right) y + \gamma  y^{3} =0
\end{equation}

Eq. \ref{eq8} is a second-order nonlinear autonomous differential equation. It can be rewritten in the form:
\begin{equation} \label{eq9}
    y_{\mathrm{ww}} + \xi_{1} y + \xi_{2} y^3 = 0,
\end{equation}
where
\begin{equation} \label{eq10}
    \begin{cases}
    \displaystyle \xi_{1} = 
        \frac{\omega^{2}}{c_{0}^{2}}
        -\frac{2 \omega}{c_{0}^{3} \beta_{2}}
        +\frac{2 k}{c_{0}^{2} \beta_{2}}
        -\frac{\Delta \beta}{c_{0}^{2} \beta_{2}},\\
    \displaystyle \xi_{2} = -\frac{2 \gamma}{c_{0}^{2} \beta_{2}}.
    \end{cases}
\end{equation}

The general solution of Eq. \ref{eq9} has pole of 1st order since two leading terms have orders \(p-2\) and \(3p\). After multiplying Eq. \ref{eq9} on \(y_{\mathrm{w}}\) and calculating integral, we obtain the equation:
\begin{equation} \label{eq11}
    \frac{y^2_{\mathrm{w}}}{2} + \frac{\xi_{1}}{2} y^2 + \frac{\xi_{2}}{4} y^4 = I_{1}.
\end{equation}

\subsection{Elliptic Legendre integral solution}
From Eq. \ref{eq11} we can obtain the wave envelope first derivative in form:
\begin{equation} \label{eq12}
    y_{\mathrm{w}} = \pm \sqrt{2I_{1} - \xi_{1}y^2-\frac{\xi_{2}}{2}y^4},
\end{equation}
or
\begin{equation} \label{eq13}
    \pm \left(\mathrm{w} - \mathrm{w}_{0}\right) = \int \frac{dy}{ \sqrt{2I_{1} - \xi_{1}y^2-\frac{\xi_{2}}{2}y^4}}.
\end{equation}

At following parameter restrictions:
\begin{equation} \label{eq13_1}
    2 \xi_{1} + \xi_{2} = 4I_{1},
\end{equation}
Integral \ref{eq13} takes the form:
\begin{equation} \label{eq13_2}
    \pm \left(\mathrm{w} - \mathrm{w}_{0}\right) = \sqrt{2I_{1}}\int_{0}^{y} \frac{dp}{ \sqrt{\left(1 - p^2\right)\left(1 - d^2 p^2\right)}},
\end{equation}
where \(d^2 = \frac{\xi_{1}}{2I_{1}} -1 =-\frac{\xi_{2}}{4I_{1}}\).

The integral in Eq. \ref{eq13_2} is a normal elliptic Legendre integral of the first kind in Jacobi form. Hence, the expression Eq. \ref{eq13}, taking into account Eq. \ref{eq13_1} lead to the following:
\begin{equation} \label{eq13_3}
    \pm \left(\mathrm{w} - \mathrm{w}_{0}\right) = \sqrt{2I_{1}} \cdot F \left(y, \sqrt{- \frac{\xi_{2}}{4I_{1}}} \right),
\end{equation}
and the integral cannot be computed analytically.

\subsection{Dark soliton solution}
In order to find exact solution of Eq. \ref{eq11}, we introduce new variable \(y\left(\mathrm{w}\right) = \sqrt{Y\left(\mathrm{w}\right)}\), and rewrite equation with lower nonlinearity degree:
\begin{equation} \label{eq14}
    Y^2_{\mathrm{w}} + 2\xi_{2}Y^3 + 4\xi_{1}Y^2 -8 I_{1} Y = 0.
\end{equation}

Next, we apply the Kudryashov Method [cite] to Eq. \ref{eq14}.
First, we define the pole order of the general solution by substituting \(Y\left(\mathrm{w}\right) = \mathrm{w}^{p}\). Leading terms have powers \(2p-2\) and \(3p\). The solution have the pole of second order.
Therefore we will be looking for \(Y\left(\mathrm{w}\right)\) in form:
\begin{equation} \label{eq15}
    Y\left(\mathrm{w}\right) = M_{0} + M_{1}Q\left(\mathrm{w}\right) + M_{2}Q^2\left(\mathrm{w}\right),
\end{equation}
where \(Q\left(\mathrm{w}\right)\) is the solution of Riccati equation:
\begin{equation}\label{eq16}
    Q_{\mathrm{w}} = \mu\,(Q^{2}-Q),
\end{equation}
of the form:
\begin{equation}\label{eq17}
    Q(\mathrm{w}) = \frac{1}{1 + \exp\left(\mu(\mathrm{w}-\mathrm{w}_{0})\right)}.
\end{equation}

Using \ref{eq15} and \ref{eq16}, Eq. \ref{eq14} can be written as:
\begin{equation}\label{eq18}
    \begin{split}
        &\left(4 \mu^{2} M_{2}^{2}+2 M_{2}^{3} \xi_{2}\right) Q^{6}+\\
        &+\left(4 \mu^{2} M_{1} M_{2}-8 \mu^{2} M_{2}^{2}+6 M_{1} M_{2}^{2} \xi_{2}\right) Q^{5}+\\
        &+\left(\mu^{2} M_{1}^{2}-8 \mu^{2} M_{1} M_{2}+4 \mu^{2} M_{2}^{2}+6 M_{0} M_{2}^{2} \xi_{2}+6 M_{1}^{2} M_{2} \xi_{2}+4 M_{2}^{2} \xi_{1}\right) Q^{4}+\\
        &+\left(-2 \mu^{2} M_{1}^{2}+4 \mu^{2} M_{1} M_{2}+12 M_{0} M_{1} M_{2} \xi_{2}+2 M_{1}^{3} \xi_{2}+8 M_{1} M_{2} \xi_{1}\right) Q^{3}+\\
        &+\left(\mu^{2} M_{1}^{2}+6 M_{0}^{2} M_{2} \xi_{2}+6 M_{0} M_{1}^{2} \xi_{2}+8 M_{0} M_{2} \xi_{1}+4 M_{1}^{2} \xi_{1}-8 \mathrm{I}_{1} M_{2}\right) Q^{2}+\\
        &+\left(6 M_{0}^{2} M_{1} \xi_{2}+8 M_{0} M_{1} \xi_{1}-8 \mathrm{I}_{1} M_{1}\right) Q+\\
        &+\left(2 \xi_{2} M_{0}^{3}+4 \xi_{1} M_{0}^{2}-8 \mathrm{I}_{1} M_{0}\right)=0.
    \end{split}
\end{equation}

Since \(Q(\mathrm{w}) \neq 0\), every coefficient of Eq. \ref{eq18} must equal zero. This lead to the following restrictions:
\begin{equation} \label{eq19}
    \begin{cases}
        \displaystyle M_{0} = \frac{8 \mathrm{I}_{1}}{\mu^{2}}\\
        \displaystyle M_{1} = -\frac{32 \mathrm{I}_{1}}{\mu^{2}}\\
        \displaystyle M_{2} = \frac{32 \mathrm{I}_{1}}{\mu^{2}}\\
        \displaystyle \xi_{1} = \frac{\mu^{2}}{2}\\
        \displaystyle \xi_{2} = -\frac{\mu^{4}}{16 \mathrm{I}_{1}}.
    \end{cases}
\end{equation}

Now the solution of an Eq. \ref{eq14} can be written as:
\begin{equation} \label{eq20}
    Y(\mathrm{w}) = \frac{8 \mathrm{I}_{1}}{\mu^{2}} - \frac{32 \mathrm{I}_{1}}{\mu^{2}} \cdot \frac{1}{1 + \exp\left(\mu(\mathrm{w}-\mathrm{w}_{0})\right)} + \frac{32 \mathrm{I}_{1}}{\mu^{2}} \cdot \frac{1}{\left(1 + \exp\left(\mu(\mathrm{w}-\mathrm{w}_{0})\right)\right)^{2}},
\end{equation}
or
\begin{equation} \label{eq21}
    Y(\mathrm{w}) = \frac{8\mathrm{I}_{1}}{\mu^2} \cdot \tanh^{2}\left(\frac{\mu\left(\mathrm{w}-\mathrm{w}_{0}\right)}{2}\right)
\end{equation}

The analytical solution of Eq. \ref{eq2} takes the form:
\begin{equation} \label{eq22}
    u(z,t) = \frac{2 \sqrt{2 I_{1}}}{\left|\mu\right|} \left|\tanh\left(\frac{\mu\left(z - c_{0} t-\mathrm{w}_{0}\right)}{2}\right)\right| \cdot \mathrm{e}^{i\left(k z-\omega t\right)}
\end{equation}

\begin{figure}
    \centering
    \includegraphics[width=0.7\linewidth]{fig1.png}
    \caption{Dark soliton (\ref{eq22}) at parameters \(I_{1} = 5.3\), \(\mu = 1.2\), \(c_{0} = 5.5\), \(t = 0\), \(\mathrm{w}_{0} = 3\)}
    \label{fig1}
\end{figure}
\subsection{Jacobi hyperbolic sine solution}
Let us consider another exact solution of Eq. \ref{eq9}, which is periodic in \(\mathrm{w}\) and is expressed through the Jacobi elliptic function \(\mathrm{sn}\). Let \(\mathrm{sn}(\cdot\,|\,m)\) be the Jacobi sine amplitude with parameter \(m \in (0,1)\). It satisfies:
\begin{equation} \label{eq11_1}
    \mathrm{sn}'^{\,2} = \left(1-\mathrm{sn}^{2}\right)\left(1-m\,\mathrm{sn}^{2}\right), \qquad \mathrm{sn}'' = -(1+m)\,\mathrm{sn} + 2m\,\mathrm{sn}^{3}.
\end{equation}

We look for the solution in form:
\begin{equation} \label{eq11_2}
    y\left(\mathrm{w}\right) = A\,\mathrm{sn}\left(\kappa\,\mathrm{w}\,|\,m\right).
\end{equation}
Then \(y_{\mathrm{ww}} = A\kappa^{2}\left(-(1+m)\,\mathrm{sn} + 2m\,\mathrm{sn}^{3}\right)\). Substituting into Eq. \ref{eq9} and collecting the coefficients of \(\mathrm{sn}\) and \(\mathrm{sn}^{3}\), we obtain:
\begin{equation} \label{eq11_3}
    \begin{cases}
        \displaystyle -\kappa^{2}\left(1+m\right) + \xi_{1} = 0,\\
        \displaystyle 2m\kappa^{2} + \xi_{2}A^{2} = 0.
    \end{cases}
\end{equation}
Hence
\begin{equation} \label{eq11_4}
    \kappa^{2} = \frac{\xi_{1}}{1+m}, \qquad A^{2} = \frac{2m\,\xi_{1}}{\left(1+m\right)\left|\xi_{2}\right|}.
\end{equation}

The period of \(\mathrm{sn}\left(\kappa\,\mathrm{w}\right)\) in \(\mathrm{w}\) is \(4K(m)/\kappa\), where \(K(m)\) is the complete elliptic integral of the first kind. The time window has length \(L\), so at fixed \(z\) the shift \(T \to T+L\) moves \(\mathrm{w}\) by \(c_{0}L\). The solution is periodic in \(T\) when
\begin{equation} \label{eq11_5}
    \frac{4K(m)}{\kappa} = c_{0}L.
\end{equation}
Since \(\kappa\) depends on \(m\) through Eq. \ref{eq11_4}, Eq. \ref{eq11_5} determines \(m\). The left side equals \(2\pi/\sqrt{\xi_{1}}\) at \(m=0\) and grows to infinity as \(m \to 1\), so a solution exists for \(c_{0}L > 2\pi/\sqrt{\xi_{1}}\).

For large \(c_{0}L\) the modulus is extremely close to 1, so \(1-m\) cannot be stored in floating point. We use the complementary modulus \(k'=\sqrt{1-m}\) instead. For the parameters of the numerical example (\texttt{run\_0.jl}) \(k' \approx 7\cdot 10^{-15}\).

The arithmetic-geometric mean (AGM) gives \(K\) and \(\mathrm{sn}\):
\begin{equation} \label{eq11_6}
    a_{0}=1,\quad b_{0}=k',\quad \tilde{c}_{0}=\sqrt{m},\qquad a_{n+1}=\frac{a_{n}+b_{n}}{2},\quad b_{n+1}=\sqrt{a_{n}b_{n}},\quad \tilde{c}_{n+1}=\frac{a_{n}-b_{n}}{2}.
\end{equation}
The iteration is continued until \(\tilde{c}_{N} \approx 0\), and then
\begin{equation} \label{eq11_7}
    K(m) = \frac{\pi}{2a_{N}}.
\end{equation}
To evaluate \(\mathrm{sn}\), we set \(\varphi_{N} = 2^{N}a_{N}u\) and iterate downward:
\begin{equation} \label{eq11_8}
    \varphi_{n-1} = \frac{1}{2}\left(\varphi_{n} + \arcsin\left(\frac{\tilde{c}_{n}}{a_{n}}\sin\varphi_{n}\right)\right), \qquad n = N,\dots,1,
\end{equation}
so that \(\mathrm{sn}(u\,|\,m) = \sin\varphi_{0}\). The modulus is found by bisection on \(\log k'\) so that Eq. \ref{eq11_5} holds.

Finally, the solution of Eq. \ref{eq2} takes the form:
\begin{equation} \label{eq11_9}
    u(z,t) = A\,\mathrm{sn}\left(\kappa\left(z - c_{0}t - \mathrm{w}_{0}\right)\,|\,m\right)\mathrm{e}^{i\left(kz-\omega t\right)}.
\end{equation}
The zeros of \(\mathrm{sn}\) are at \(\mathrm{w}=0\) and \(\mathrm{w}=2K/\kappa\). Each zero is a dark soliton with a \(\pi\) phase jump, and neighbouring zeros have opposite polarity. Within one window there are two such solitons, half a window apart.

For \(m \to 1\) we have \(\mathrm{sn} \to \tanh\), \(\kappa \to \sqrt{\xi_{1}/2}\) and \(A \to \sqrt{\xi_{1}/\left|\xi_{2}\right|}\). With \(\xi_{1} = \mu^{2}/2\) and \(\xi_{2} = -\mu^{4}/(16 I_{1})\) from Eq. \ref{eq19}, this gives \(\kappa = \mu/2\) and \(A = 2\sqrt{2I_{1}}/|\mu|\). These are the parameters of the single dark soliton Eq. \ref{eq22}, which is not periodic.

\section{Numerical Analysis}
    To build a numerical scheme for numerical analysis, we rewrite original system (\ref{eq1}) in the form:
    \begin{equation}
        \frac{\partial}{\partial z}\begin{pmatrix}u\\v\end{pmatrix}
        = \bigl(\hat L + \hat N\bigr)\begin{pmatrix}u\\v\end{pmatrix},
    \end{equation}
    where
    \begin{equation}
        \hat L =
            \left(\frac{i\Delta\beta}{2}+\frac{g}{2}\right)
            - \delta\frac{\partial}{\partial t}
            + \left(\frac{g}{2\Omega_{g}^2}-\frac{i\beta_{2}}{2}\right)\frac{\partial^2}{\partial t^2}
    \end{equation}
    contains all linear terms (dispersion, birefringence, and frequency‑dependent gain) and
    \begin{equation}\label{eq:nonlin_ode}
        \hat N\begin{pmatrix}u\\v\end{pmatrix} =
            i\gamma \cdot \begin{pmatrix}
            \Bigl[\bigl(|u|^2+\tfrac{2}{3}|v|^2\bigr)u + \tfrac{1}{3}v^2 u^*\Bigr]\\[4pt]
            \Bigl[\bigl(|v|^2+\tfrac{2}{3}|u|^2\bigr)v + \tfrac{1}{3}u^2 v^*\Bigr]
            \end{pmatrix}
    \end{equation}
    is the nonlinear term.

    We employ a split‑step Fourier method with symmetric (Strang) splitting and formulate initial-boundary problem with periodic boundary conditions:
    \begin{equation} 
        \begin{cases}
            \displaystyle u\left(z, -\frac{L_{t}}{2}\right) = u\left(z, \frac{L_{t}}{2}\right),\\
            \displaystyle v\left(z, -\frac{L_{t}}{2}\right) = v\left(z, \frac{L_{t}}{2}\right),\\
            \displaystyle \frac{\partial u}{\partial t}\left(z, -\frac{L_{t}}{2}\right) = \cfrac{\partial u}{\partial t}\left(z, \frac{L_{t}}{2}\right),\\
            \displaystyle \frac{\partial v}{\partial t}\left(z, -\frac{L_{t}}{2}\right) = \cfrac{\partial v}{\partial t}\left(z, \frac{L_{t}}{2}\right).
        \end{cases}
    \end{equation}
    The propagation is performed on a periodic time window $t\in[\frac{-L_{t}}{2},\frac{L_{t}}{2}]$ sampled by $N_t$ equally spaced points $\Delta t = L_t/N_t$.

    The solution over a step $\Delta z$ is approximated by the symmetric splitting
    \begin{equation}
        \begin{pmatrix} u\\v \end{pmatrix}_{z+\Delta z} \approx
        e^{\frac{\Delta z}{2}\hat L} \cdot
        e^{\Delta z\hat N} \cdot
        e^{\frac{\Delta z}{2}\hat L}
        \begin{pmatrix}u\\v\end{pmatrix}_{z}.
    \end{equation}
    The linear exponential $e^{\frac{\Delta z}{2}\hat L}$ is applied in the Fourier domain,
    while the nonlinear exponential $e^{\Delta z\hat N}$ is integrated with a fourth‑order Runge–Kutta (RK4) scheme.
    
    \subsection{Linear sub‑step}
    In the Fourier domain the linear operator becomes a multiplicative factor. With the convention $\partial/\partial t \leftrightarrow i\omega_{f}$,
    the Fourier transform of $\hat L$ acting on a single field $u$ (or $v$) is the following:
    \begin{equation}
        \widetilde L(\omega_{f})
        = \frac{g}{2} - \frac{g}{2\Omega_g^2}\omega_{f}^2
        + i\Bigl(\frac{\Delta\beta}{2} - \delta\omega_{f} + \frac{\beta_2}{2}\omega_{f}^2\Bigr).
    \end{equation}
    Both fields $u$ and $v$ experience the same linear response because the linear part
    of Eq. \ref{eq1} is diagonal and identical for the two polarisations.
    Hence a half‑step of the linear operator is applied in the frequency domain as
    \begin{equation}
        \left[
        \begin{array}{ll}
            \tilde u\left(\omega_f\right) \Big| _{z+\frac{\Delta z}{2}} =
            \exp\!\Bigl[\widetilde L(\omega_f)\frac{\Delta z}{2}\Bigr] \tilde u\left(\omega_f\right) \Big| _{z},\\
            \tilde v\left(\omega_f\right) \Big| _{z+\frac{\Delta z}{2}} =
            \exp\!\Bigl[\widetilde L(\omega_f)\frac{\Delta z}{2}\Bigr] \tilde v\left(\omega_f\right) \Big| _{z},
        \end{array}
        \right.
    \end{equation}
    where $\tilde u = \mathcal{F}[u]$ and $\tilde v = \mathcal{F}[v]$ are obtained by the fast Fourier transform (FFT).
    
    After the nonlinear full-step, the linear half-step applied again. Then, the fields are transformed back to the time domain with the inverse FFT.
    The FFT routines assume periodic boundary conditions, consistent with the finite sampling window.
    
    \subsection{Nonlinear sub‑step}
    The nonlinear sub‑step advances the ordinary differential equation
    \begin{equation} \label{nonlinear_eq}
        \frac{\partial}{\partial z}\begin{pmatrix}u\\v\end{pmatrix}
        = \hat N\begin{pmatrix}u\\v\end{pmatrix},
    \end{equation}
    over a full step $\Delta z$.  We solve \eqref{nonlinear_eq} with the classical fourth‑order Runge–Kutta method.
    For given nonlinear term, the integration can be performed independently
    at each grid point using vectorized operations since there is no derivatives and neighbor-dependencies. We use classic RK4 in the form:
    \begin{align}
        &\mathbf{k}_1 = \hat N(\mathbf{w}), \nonumber\\
        &\mathbf{k}_2 = \hat N(\mathbf{w} + \tfrac{\Delta z}{2}\mathbf{k}_1), \nonumber\\
        &\mathbf{k}_3 = \hat N(\mathbf{w} + \tfrac{\Delta z}{2}\mathbf{k}_2), \nonumber\\
        &\mathbf{k}_4 = \hat N(\mathbf{w} + \Delta z\,\mathbf{k}_3), \nonumber\\
        &\mathbf{w}\Big| _{z+\Delta z} = \mathbf{w}\Big| _{z} + \frac{\Delta z}{6}\bigl(\mathbf{k}_1+2\mathbf{k}_2+2\mathbf{k}_3+\mathbf{k}_4\bigr), \nonumber
    \end{align}
    where $\mathbf{w}=(u,v)^T$. The approximation order RK4 ensures that the overall splitting error remains the dominant source of numerical inaccuracy.
    
    \subsection{Symmetric split‑step scheme}
    Combining the two half‑steps, one complete step from $z$ to $z+\Delta z$ proceeds as follows:
    \begin{enumerate}
        \item \textbf{First linear half-step.} Compute $\tilde u = \mathcal{F}[u]$, $\tilde v = \mathcal{F}[v]$, by FFT.
        Multiply by $\exp[\widetilde L(\omega)\frac{\Delta z}{2}]$, and transform back:
        $u \leftarrow \mathcal{F}^{-1}[\tilde u]$, $v \leftarrow \mathcal{F}^{-1}[\tilde v]$.
        \item \textbf{Full nonlinear step.} Advance $(u,v)$ by $\Delta z$ using the RK4 integration of \eqref{eq:nonlin_ode}.
        \item \textbf{Second linear half-step.} Repeat step~1.
    \end{enumerate}
    
\section{Modeling Results}
    We consider dark soliton of the form (\ref{eq21}) as initial conditions for initial-boundary problem for system (\ref{eq1}). We take \(v(z,t) \equiv 0\). Fig. \ref{fig2} illustrates the pulse behavior under certain model parameters. The soliton propagates steadily.
    \begin{figure}
        \centering
            \begin{subfigure}[t]{0.49\linewidth}
                \centering
                \includegraphics[width=\linewidth]{fig2.png}
                \subcaption{Numerical solution \(u\) profile at \(z=0\) and \(z=100\).}
            \end{subfigure}
            \hfill
            \begin{subfigure}[t]{0.49\linewidth}
                \centering
                \includegraphics[width=\linewidth]{fig3.png}
                \subcaption{\(\left|u(z,t)\right|\)}
            \end{subfigure}
        \caption{
            \\Pulse parameters:
            \(k = 12.6\),
            \(c_{0} = 1.4\),
            \(\mu = 0.97\),
            \(I_{1} = 1\),
            \(\omega = 0\),
            \(\mathrm{w}_{0} = 15\).\\
            Model parameters:
            \(\beta_{2} = 0.025\),
            \(\Omega_{g} = 10\),
            \(\delta = 20\),
            \(\gamma = 0.0013\),
            \(g = 0\),
            \(\Delta \beta = 5\).
        }
        \label{fig2}
    \end{figure}

    By increasing the coupling coefficient \(\gamma\) and perturbing initial condition \(v(z=0, t) = \sin(t)\cdot\exp\left(-(c_{0}t-w_{0})\right)\), we can observe pulse decay, as depicted on Fig. \ref{fig3}.
    \begin{figure}
        \centering
            \begin{subfigure}[t]{0.49\linewidth}
                \centering
                \includegraphics[width=\linewidth]{fig4.png}
                \subcaption{\(\left|u(z,t)\right|\)}
            \end{subfigure}
            \hfill
            \begin{subfigure}[t]{0.49\linewidth}
                \centering
                \includegraphics[width=\linewidth]{fig5.png}
                \subcaption{\(\left|v(z,t)\right|\)}
            \end{subfigure}
        \caption{
            Numerical solution for disturbed \(v(z=0, t)\). \\
            Model parameters:
            \(\beta_{2} = 0.025\),
            \(\Omega_{g} = 10\),
            \(\delta = 20\),
            \(\gamma = 0.13\),
            \(g = 0\),
            \(\Delta \beta = 5\).
        }
        \label{fig3}
    \end{figure}

    From numerical analysis, the following conclusions are made:
    \begin{enumerate}
        \item Dark soliton, found by Kudryashov method, propagate steadily;
        \item Model parameter \(\delta\) is a linear-transfer coefficient;
        \item Model parameter \(\gamma\) is a coupling parameter;
        \item Model parameter \(g\) is a linear amplification parameter;
    \end{enumerate}
\end{document}

